% GENERATED FILE. Edit canonical lessons, then run npm run generate:books.
% canonical-source-hash: fnv1a64:e3b4ae3cf9ff9c8d
% canonical-lessons: PT-C509-baliza, PT-C509-penalti, PT-C509-avancado, PT-C509-rematar, PT-C509-cabecear

\chapter{Na baliza}
\label{ch:pt-na-baliza}
\hlchaptermodality{509}

\begin{chapteropening}
\textbf{By the end of this chapter:} \emph{I can talk about goals, goalkeepers and strikers.}

Complete the last lesson of chapter 509: I can talk about goals, goalkeepers and strikers.
\end{chapteropening}

\section[a baliza]{a baliza — goal (the posts and net)}

\label{lesson:PT-C509-baliza}



\begin{quote}
Before the new one: say the Portuguese for stand, then the Portuguese for substitute.
\end{quote}

\subsection*{You'll want to know: a baliza}
\textbf{a baliza} — \textquotedblleft{}goal (the posts and net)\textquotedblright{}.

\begin{quote}
A bola bateu no poste e não entrou na baliza.
\end{quote}

\begin{quote}
\emph{The ball hit the post and did not go into the goal.}
\end{quote}

In Brazil: \emph{o gol}. A scored goal is \emph{um golo} in Portugal.

\begin{cousinweb}
The origin of \emph{baliza}, once a marker post or a buoy, is not settled.
\end{cousinweb}

\begin{grammarlens}[title={What you've built}]
One more word for telling, explaining and describing at B1.
\end{grammarlens}

\subsection*{Your turn}
\begin{itemize}
\raggedright
  \item \emph{Say it:} \emph{a baliza}
  \item \emph{Say it:} \emph{a baliza}, once more
  \item \emph{Say it:} a sentence of your own with \emph{a baliza}
  \item \emph{Recall:} say the Portuguese for stand, then the Portuguese for substitute, then say \emph{a baliza} again
\end{itemize}

\begin{tcolorbox}[breakable,colback=teal!4,colframe=teal!35!black,title={Before you move on}]
What does a baliza mean? (\textquotedblleft{}Goal (the posts and net)\textquotedblright{}.) Say it once more.
\end{tcolorbox}
\section[o penálti]{o penálti — penalty (kick)}

\label{lesson:PT-C509-penalti}



\begin{quote}
Before the new one: say the Portuguese for substitute, then the Portuguese for goal.
\end{quote}

\subsection*{You'll want to know: o penálti}
\textbf{o penálti} — \textquotedblleft{}penalty (kick)\textquotedblright{}.

\begin{quote}
O árbitro marcou um penálti no último minuto.
\end{quote}

\begin{quote}
\emph{The referee gave a penalty in the last minute.}
\end{quote}

In Brazil: \emph{o pênalti}. A more formal name is \emph{a grande penalidade}.

\begin{cousinweb}
\emph{penálti} was borrowed from English \emph{penalty}, which goes back through French and Medieval Latin to Latin \emph{poena} \textquotedblleft{}a punishment\textquotedblright{} (Lewis \& Short).
\end{cousinweb}

\begin{grammarlens}[title={What you've built}]
One more word for telling, explaining and describing at B1.
\end{grammarlens}

\subsection*{Your turn}
\begin{itemize}
\raggedright
  \item \emph{Say it:} \emph{o penálti}
  \item \emph{Say it:} \emph{o penálti}, once more
  \item \emph{Say it:} a sentence of your own with \emph{o penálti}
  \item \emph{Recall:} say the Portuguese for substitute, then the Portuguese for goal, then say \emph{o penálti} again
\end{itemize}

\begin{tcolorbox}[breakable,colback=teal!4,colframe=teal!35!black,title={Before you move on}]
What does o penálti mean? (\textquotedblleft{}Penalty (kick)\textquotedblright{}.) Say it once more.
\end{tcolorbox}
\section[o avançado]{o avançado — forward, striker}

\label{lesson:PT-C509-avancado}



\begin{quote}
Before the new one: say the Portuguese for goal, then the Portuguese for penalty.
\end{quote}

\subsection*{You'll want to know: o avançado}
\textbf{o avançado} — \textquotedblleft{}forward, striker\textquotedblright{}.

\begin{quote}
O avançado marcou dois golos.
\end{quote}

\begin{quote}
\emph{The striker scored two goals.}
\end{quote}

In Brazil: \emph{o atacante}.

\begin{cousinweb}
\emph{avançado} is built on \emph{avançar} \textquotedblleft{}to move forward\textquotedblright{}, which goes back to Vulgar Latin \emph{abanteāre} (EtymDB), built on \emph{ab ante} \textquotedblleft{}from in front\textquotedblright{}.
\end{cousinweb}

\begin{grammarlens}[title={What you've built}]
One more word for telling, explaining and describing at B1.
\end{grammarlens}

\subsection*{Your turn}
\begin{itemize}
\raggedright
  \item \emph{Say it:} \emph{o avançado}
  \item \emph{Say it:} \emph{o avançado}, once more
  \item \emph{Say it:} a sentence of your own with \emph{o avançado}
  \item \emph{Recall:} say the Portuguese for goal, then the Portuguese for penalty, then say \emph{o avançado} again
\end{itemize}

\begin{tcolorbox}[breakable,colback=teal!4,colframe=teal!35!black,title={Before you move on}]
What does o avançado mean? (\textquotedblleft{}Forward, striker\textquotedblright{}.) Say it once more.
\end{tcolorbox}
\section[rematar]{rematar — to shoot (at goal); to finish off}

\label{lesson:PT-C509-rematar}



\begin{quote}
Before the new one: say the Portuguese for penalty, then the Portuguese for forward.
\end{quote}

\subsection*{You'll want to know: rematar}
\textbf{rematar} — \textquotedblleft{}to shoot (at goal); to finish off\textquotedblright{}.

\begin{quote}
O avançado rematou com o pé esquerdo.
\end{quote}

\begin{quote}
\emph{The striker shot with his left foot.}
\end{quote}

A shot is \emph{um remate}.

\begin{cousinweb}
\emph{rematar} is built from \emph{re-} and \emph{matar} in its old sense \textquotedblleft{}to finish, to end\textquotedblright{}; \emph{matar} is usually traced to Latin \emph{mactāre} \textquotedblleft{}to slaughter\textquotedblright{} (EtymDB).
\end{cousinweb}

\begin{grammarlens}[title={What you've built}]
One more word for telling, explaining and describing at B1.
\end{grammarlens}

\subsection*{Your turn}
\begin{itemize}
\raggedright
  \item \emph{Say it:} \emph{rematar}
  \item \emph{Say it:} \emph{rematar}, once more
  \item \emph{Say it:} a sentence of your own with \emph{rematar}
  \item \emph{Recall:} say the Portuguese for penalty, then the Portuguese for forward, then say \emph{rematar} again
\end{itemize}

\begin{tcolorbox}[breakable,colback=teal!4,colframe=teal!35!black,title={Before you move on}]
What does rematar mean? (\textquotedblleft{}To shoot (at goal); to finish off\textquotedblright{}.) Say it once more.
\end{tcolorbox}
\section[cabecear]{cabecear — to head (the ball)}

\label{lesson:PT-C509-cabecear}



\begin{quote}
Before the new one: say the Portuguese for forward, then the Portuguese for to shoot.
\end{quote}

\subsection*{You'll want to know: cabecear}
\textbf{cabecear} — \textquotedblleft{}to head (the ball)\textquotedblright{}.

\begin{quote}
O defesa cabeceou a bola para fora.
\end{quote}

\begin{quote}
\emph{The defender headed the ball out.}
\end{quote}

Built on \emph{a cabeça}.

\begin{cousinweb}
\emph{cabecear} is built on \emph{cabeça} \textquotedblleft{}head\textquotedblright{}, which goes back to a Vulgar Latin form of Latin \emph{caput} \textquotedblleft{}head\textquotedblright{} (Lewis \& Short).
\end{cousinweb}

\begin{grammarlens}[title={What you've built}]
One more word for telling, explaining and describing at B1.
\end{grammarlens}

\subsection*{Your turn}
\begin{itemize}
\raggedright
  \item \emph{Say it:} \emph{cabecear}
  \item \emph{Say it:} \emph{cabecear}, once more
  \item \emph{Say it:} a sentence of your own with \emph{cabecear}
  \item \emph{Recall:} say the Portuguese for forward, then the Portuguese for to shoot, then say \emph{cabecear} again
\end{itemize}

\begin{tcolorbox}[breakable,colback=teal!4,colframe=teal!35!black,title={Before you move on}]
What does cabecear mean? (\textquotedblleft{}To head (the ball)\textquotedblright{}.) Say it once more.
\end{tcolorbox}
